% 20-minute viva voce deck — 14 content slides + backups.
% Pacing target: ~85s/slide; script in viva_script.md.
\documentclass[11pt,aspectratio=169]{beamer}
\usetheme{metropolis}
\usepackage{booktabs}
\newcommand{\ZVF}{\ensuremath{\mathrm{ZVF}}}
\setbeamertemplate{frame numbering}[counter]

\title{The Zero-Variance Fraction}
\subtitle{Diagnosing and Budgeting Signal Starvation in Group-Relative RL Post-Training}
\author{Arvind C R \\ \small Guide: Prof.\ Ramesh Prakash Guledgudd}
\institute{M.Tech Viva Voce \, · \, PES University}
\date{July 2026}

\begin{document}
\maketitle

% ---- 1 ----
\begin{frame}{The problem in one picture}
\begin{columns}
\column{0.55\textwidth}
GRPO trains on \textbf{groups} of $G$ completions per prompt; advantages are
centred within the group.\\[6pt]
\textbf{If every completion in a group gets the same reward, the group
contributes \alert{zero gradient}.}\\[6pt]
Two ways it happens:
\begin{itemize}
  \item all wrong — task too hard
  \item all right — task \emph{mastered}
\end{itemize}
\column{0.42\textwidth}
\begin{block}{The blind spot}
The second case is invisible in the reward curve: reward reads
\emph{success} exactly when learning has \emph{stopped}.
\end{block}
\end{columns}
\end{frame}

% ---- 2 ----
\begin{frame}{The instrument}
\begin{definition}
$\displaystyle \ZVF_t=\frac{1}{|\mathcal{P}_t|}\sum_{p}
\mathbf{1}\big[\mathrm{Var}\,r=0 \text{ within the group}\big]$
\qquad (gradient utilisation $=1-\ZVF_t$)
\end{definition}
\begin{itemize}
  \item Costs nothing: computed from rewards already in hand
  \item Population form under binary rewards: $h_G(p)=p^G+(1-p)^G$ —
        a \textbf{U-shape} in difficulty
  \item Must be read \textbf{with reward}: \ZVF{} alone aliases mastery
        with incapacity
\end{itemize}
\end{frame}

% ---- 3 ----
\begin{frame}{Claim contract (what I defend — and what I don't)}
\textbf{Claim 1.} \ZVF{}+reward is a cheap, calibrated online diagnostic of
zero-advantage regimes.\\[4pt]
\textbf{Claim 2.} At matched rollout budget, group size selects \emph{which
end} of training starves ($G{=}2$: all-correct wall; $G{=}16$: none).\\[8pt]
\begin{alertblock}{Explicit non-claims}
no predictive/causal power for \ZVF{} · no data-dependent optimal $G$ ·
no controller-efficacy claim · no generalisation beyond
Qwen3-8B / GSM8K-family / one managed LoRA stack, 1--3 seeds
\end{alertblock}
\end{frame}

% ---- 4 ----
\begin{frame}{Scope and infrastructure}
\begin{itemize}
  \item One 8B model (Qwen3-8B, LoRA r4) on a managed training API;
        vLLM eval on 3 GPU backends; 7-library observational corpus
        (0.6B–671B logged runs)
  \item Base capability, GSM8K test ($n{=}32$/problem, clustered 95\% CIs):
        \textbf{pass@1 30.4\% — pass@32 91.0\%} $\Rightarrow$ little
        capability headroom; claims scoped accordingly
  \item Every result: seeded, manifest-ed, W\&B-mirrored, resumable
        checkpoints; partial runs stamped, never promoted
\end{itemize}
\end{frame}

% ---- 5 ----
\begin{frame}{Claim 1 — the all-correct wall (the key exhibit)}
Matched budget: \textbf{2{,}560 rollouts per arm}, two seeds each\\
$G{=}2\times160$ steps \quad vs. \quad $G{=}16\times20$ steps
\begin{center}
\begin{tabular}{@{}lcc@{}}
\toprule
 & $G{=}2$ & $G{=}16$\\
\midrule
final train reward & $\approx 1.0$ & comparable\\
late-run \ZVF{} & \alert{$0.75$–$1.0$} & $0.00$–$0.25$\\
gradient signal at end & \alert{none} & sustained\\
\bottomrule
\end{tabular}
\end{center}
\medskip
Reward says ``succeeding''; \ZVF{} says ``training is over.''
\textbf{Only the pair tells the truth.}
\end{frame}

% ---- 6 ----
\begin{frame}{Theory 1/3 — calibration (T1)}
\begin{itemize}
  \item $\ZVF_t$ is an unbiased binomial-proportion estimator; we ship it
        with a \textbf{Wilson interval — empirical coverage 0.95–0.98 in
        every tested setting}
  \item Curriculum-ordered batches collapse \emph{global} coverage
        (0.08–0.40) but the interval stays calibrated for the
        \emph{local} estimand (0.944) — a labelling requirement,
        not an invalidation
\end{itemize}
\end{frame}

% ---- 7 ----
\begin{frame}{Theory 2/3 — the reliability budget (T2)}
\[
N(\ZVF)=G\Big\lceil \tfrac{\ln\delta}{\ln \ZVF}\Big\rceil
\quad = \quad
\text{the } (1{-}\delta)\text{-quantile of rollouts to the next
informative group}
\]
\begin{itemize}
  \item A \textbf{budget}, not a minimum — signal can arrive in the first
        $G$ rollouts w.p.\ $1-\ZVF$
  \item Empirically \textbf{exact}: model quantile $=$ observed quantile,
        ratio 1.00, in all six difficulty strata
  \item Hardest stratum ($\hat p\!\approx\!0.01$): 160 rollouts to
        \emph{guarantee} one usable gradient at $\delta{=}0.1$
\end{itemize}
\end{frame}

% ---- 8 ----
\begin{frame}{Theory 3/3 — an honest negative result (T3)}
Signal-per-rollout objective:
$J(G)=\frac1G\,\mathbb{E}_\phi[(1-h_G(p))\,p(1-p)]$
\begin{itemize}
  \item Algebraic identity: $1-p^G-(1-p)^G = G\,p(1-p)$ at $G\in\{2,3\}$
  \item $\Rightarrow\; J(2)=J(3)=\mathbb{E}[(p(1-p))^2]$
        \textbf{for every prior} — argmax always $\{2,3\}$
  \item Our earlier ``theory--experiment agreement'' was therefore
        \alert{guaranteed, not a test} — found in external review,
        corrected openly
  \item Consequence: per-rollout accounting cannot justify adaptive $G$;
        richer objectives required
\end{itemize}
\end{frame}

% ---- 9 ----
\begin{frame}{Claim 2 — group size is a schedule}
\begin{itemize}
  \item Static sweeps are inconclusive (non-monotone at fixed steps;
        optimum drifts to $G{\approx}32$ at fixed token budget)
  \item The matched-\textbf{rollout}-budget view resolves it:
        small $G$ = more optimiser steps early, starves late;
        large $G$ = fewer steps, signal survives the endgame
  \item Contrastive ($G{=}2\!\sim\!G{=}16$) equivalence claims hold on
        final reward, \textbf{not on signal availability} — and decay from
        97.6\% to 72.7\% retention as training scale grows
  \item Controller: designed, evidence-gapped, \textbf{explicitly not
        claimed} pending compute-matched trials vs static $G{=}16$
\end{itemize}
\end{frame}

% ---- 10 ----
\begin{frame}{Loss form — GRPO vs Dr.GRPO (six uncapped arms)}
Qwen3-8B, 1{,}024-token cap, 30 steps, $G{=}8$, \textbf{3 seeds per loss}
\begin{center}
\small
\begin{tabular}{@{}lcc@{}}
\toprule
 & length first-5$\to$last-10 & late \ZVF{}\\
\midrule
GRPO & $1004{\to}905$,\; $981{\to}944$,\; $996{\to}900$ & $0.47/0.47/0.55$\\
Dr.GRPO & $999{\to}931$,\; $972{\to}902$,\; $1000{\to}878$ & $0.45/0.72/0.70$\\
\bottomrule
\end{tabular}
\end{center}
\medskip
\textbf{No length inflation in either loss; no \ZVF{} separation.}
At this scale the loss-form choice has no observable footprint.
\end{frame}

% ---- 11 ----
\begin{frame}{The incident that became a chapter}
\begin{block}{The unwired flag (2026-07-11)}
A six-arm GRPO-vs-Dr.GRPO panel silently trained \textbf{one} objective:
the \texttt{--loss} flag was documented but never wired. No reward, length,
or \ZVF{} trace revealed it — only reading the runner did.
\end{block}
\begin{itemize}
  \item Response: invalidate loudly, preserve artifacts, fix, rerun same day
  \item The corrected panel \alert{reversed} one conclusion of the invalid one
  \item Lesson: \textbf{stack identity cannot be certified from
        plausible-looking outputs}
\end{itemize}
\end{frame}

% ---- 12 ----
\begin{frame}{From incident to standard}
\textbf{Eight-item minimum reportable stack} — every item once flipped a
result in our own corpus:
\begin{enumerate}\small
  \item loss form \; 2. reference policy/KL \; 3. sampler/backend/precision
        + \textbf{base-checkpoint identity} \; 4. per-step \ZVF{} trajectory
  \setcounter{enumi}{4}
  \item group-size schedule \; 6. held-out split \; 7. decontamination +
        parser probe \; 8. \textbf{pass@$k$ curves}, not pass@1 alone
\end{enumerate}
\medskip
Documented levers: backend swap $\Rightarrow$ \textbf{17$\times$} reward span;
same ``DAPO'' label $\Rightarrow$ \ZVF{} 0.00 vs 0.58; parser micro-jitter
$\Rightarrow$ \ZVF{} $0.158\to0.000$. Machine-readable registry +
flip-risk \texttt{stackdiff} ship with the bench.
\end{frame}

% ---- 13 ----
\begin{frame}{Limitations and roadmap}
\textbf{Limitations} — one model, one task family, one closed stack,
1–3 seeds; theory has declared proof gaps; \ZVF{} is binary-reward-specific.
\\[8pt]
\textbf{Post-degree publication plan (externally reviewed, gated):}
\begin{itemize}
  \item Diagnostic paper (Claims 1–2) — ready without controller claims
  \item Controller paper — gated on compute-matched win vs static $G{=}16$
        \emph{and} reward-only schedules, $\ge3$ seeds
  \item Survival audit — gated on an open stack + differential-objective
        test tooling
\end{itemize}
\end{frame}

% ---- 14 ----
\begin{frame}{Conclusion}
\begin{center}
\Large Under verifiable rewards, group-relative training starves at
\textbf{both ends} of difficulty.\\[10pt]
\ZVF{} + reward sees both walls — including the one the reward curve
cannot show.\\[10pt]
Group size decides \textbf{which wall} your budget hits.\\[14pt]
\normalsize Thesis, 17 consolidated working papers, and every artifact:
\texttt{github.com/arvindcr4/tinker-rl-lab}
\end{center}
\end{frame}

% ---- Backups ----
\appendix
\begin{frame}{Backup: provenance of the 17 papers}
Benchmark family (4) $\to$ Ch.3 · P2+ZVF audit $\to$ Ch.4 · theory $\to$ Ch.5
· P3+P7+E-R2b $\to$ Ch.6 · P4+E-P4 $\to$ Ch.7 · P1 $\to$ Ch.8 ·
P5/P6/position/registry/audit $\to$ Ch.9 · P8 out-of-programme.
Canonicalisation per \texttt{PAPERS\_README.md}; every claim $\to$ artifact
in thesis Appendix~B.
\end{frame}

\begin{frame}{Backup: key numbers}
\small
GSM8K base pass@1/8/32: $30.4/79.7/91.0\%$ (clustered CIs) ·
train-pool \ZVF{} $10.86\%$, all all-incorrect ·
T1 Wilson coverage $0.95$–$0.98$ ·
T2 fit ratio $1.00$ (6 strata) ·
E-R2b: $\ZVF$ $0.75$–$1.0$ vs $\le0.25$ at 2{,}560 rollouts ·
E-P4: lengths shrink 6–12\% both losses ·
MBPP transfer: zero forgetting; pass@32 frontier $+1.5$–$+2.5$ pts (1 seed) ·
MATH-500: GSM8K gains do not replicate (partial) ·
held-out GSM8K control: $82.0\to83.3\%$, $p{=}0.26$.
\end{frame}

\begin{frame}{Backup: corrections adopted from review}
T2 quantifier error (quantile $\ne$ minimum) · T2 validation direction ·
T3 universality (J(2)=J(3)) · GU definition · off-by-one final target in 7
trainers · generation seeding · parser v2 (last-boxed) · exact hypergeometric
T1 truth. All fixed 2026-07-11 (commits \texttt{5486ae2}, \texttt{acc08eb});
a further audit pass produced only fabricated findings — all refuted against
source, reinforcing the verify-against-the-stack thesis.
\end{frame}

\end{document}
